% Book adaptation of the supplied nonclosed regularity manuscript.
% Review scope: research/editorial/nonclosed-interior-regularity-review.md.
\section{Interior smoothness on the regular locus}
\label{sec:nonclosed-interior-smoothness}

Throughout this section, $\theta$ is a smooth big real $(1,1)$-form, and $\rho$, $Z=\{\rho=-\infty\}$ and $\Omega=X\setminus Z$ are the strict analytic barrier and its regular locus fixed in \cref{sec:big-minimal-comparison}. Thus $\rho\leq0$ has analytic singularities with smooth remainder, and $\theta+\ddc\rho\geq\omega$. If $\varphi\in\PSH(X,\theta)$ has minimal singularities, then $\varphi\geq V_\theta-C\geq\rho-C$ for a constant $C$. Hence $\varphi$ is locally bounded on $\Omega$, and its Monge--Amp\`ere measure there is the bounded-potential product.

For such a potential, \cref{thm:big-minimal-holder} gives local H\"older continuity on $\Omega$ when the density belongs to $L^p$ for some $p>1$. We show that a smooth strictly positive density on the whole of $\Omega$ then makes the potential smooth on $\Omega$. Local information alone cannot give this: \cref{ex:local-lipschitz-constant-measure} below is a locally Lipschitz plurisubharmonic function with constant Monge--Amp\`ere density that is not $C^1$. The argument therefore uses the barrier globally. By the proper gap of \cref{lem:interior-proper-gap}, the difference between the solution and the barrier attains a minimum on the closed set where smoothness fails, and strict positivity of $\theta+\ddc\rho$ turns this minimum into a quadratic separation. Near the minimum, smooth Dirichlet approximants have traces of polynomial growth; the separation improves this growth until a uniform trace bound yields elliptic regularity.

\Needspace{12\baselineskip}
\begin{theorem}\label{thm:nonclosed-interior-smoothness}\index{Interior regularity}
Suppose $\varphi\in\PSH(X,\theta)$ has minimal singularities and is locally
H\"older continuous on $\Omega$. Let $c>0$ and let $f$ be smooth and strictly positive on $\Omega$. Assume
\[
 (\theta+\ddc\varphi)^n=cf\,\omega^n\quad\text{on }\Omega.
\]
Then
$\varphi\in C^\infty(\Omega)$ and $\theta+\ddc\varphi>0$ there.
\end{theorem}

The H\"older hypothesis holds when, in addition, $f\in L^p(X,\omega^n)$ for some $p>1$, by \cref{thm:big-minimal-holder}; see \cref{cor:big-lp-interior-smoothness}. The theorem also applies when local H\"older continuity is obtained by another argument.

The proof is by contradiction and occupies the rest of the section. Let $\mathcal S\subset\Omega$ be the closed set of points near which $\varphi$ is not smooth, suppose that $\mathcal S\neq\varnothing$, and let $p\in\mathcal S$ minimize $\varphi-\rho$ on $\mathcal S$. On a small ball centered at $p$ the argument has four steps.
\begin{enumerate}
\item A comparison principle for continuous solutions (\cref{sec:interior-calculus}) and the smooth Dirichlet problem give smooth solutions $u_j$ of the equation that converge uniformly to $\varphi$ at a polynomial rate in $j$, while the traces of $\theta+\ddc u_j$, measured in Euclidean coordinates, grow at most polynomially in $j$ (\cref{sec:interior-approximants}).
\item Near $p$, the quadratic separation along $\mathcal S$, together with uniform trace bounds on compact sets where $\varphi$ is already smooth, improves the trace growth to $O(j^\gamma)$ for every $\gamma>0$. Interpolation with the convergence rate then gives uniform $C^{1,\alpha}$ bounds for every $\alpha<1$ (\cref{sec:interior-localization}).
\item A third-order estimate whose constants depend polynomially on the ellipticity turns uniform $C^{1,\alpha}$ bounds, for $\alpha$ close to one, into a uniform trace bound (\cref{sec:interior-calabi}).
\item A continuous solution with bounded trace is smooth (\cref{sec:interior-criterion}). Hence $\varphi$ is smooth near $p$, contradicting $p\in\mathcal S$ (\cref{sec:interior-main-proof}).
\end{enumerate}

\subsection{Local comparison}\label{sec:interior-calculus}

The comparison principle below, for continuous sub- and supersolutions on a bounded domain, identifies the given continuous potential with limits of smooth Dirichlet solutions. Since the background need not be closed, the classical Bedford--Taylor comparison does not apply directly. We freeze the background at a point, apply the closed-current comparison formula \eqref{eq:interior-power-comparison} to quadratic modifications, and control the error by the mixed inequalities of \cref{lem:interior-mixed-background,lem:interior-convolution}.

The local arguments of this section take place in coordinate balls, always with compact closure in a larger coordinate chart. There we put $\beta=\ddc|z|^2$, write $\partial_i=\partial/\partial z_i$ and $\partial_{\bar i}=\partial/\partial\bar z_i$, and take coefficient matrices of real $(1,1)$-forms relative to $\beta$. For such a form $A$, $\operatorname{tr}_\beta A$ and $\det_\beta A$ denote the trace and the determinant of its coefficient matrix, so that $A^n=(\det_\beta A)\beta^n$, and $\Delta_\beta=\operatorname{tr}_\beta\ddc$. For a positive form $T$, $\operatorname{tr}_T$ is the trace computed with the inverse of the coefficient matrix of $T$. Indices run from $1$ to $n$, repeated indices in tensor expressions are summed, and $I$ denotes the identity matrix. In local statements, $\vartheta$ denotes a smooth real $(1,1)$-form, not necessarily closed, which plays the role of $\theta$.

\begin{proposition}\label{prop:interior-comparison}
Let $D\Subset\mathbb C^n$ be a bounded domain and let $\vartheta$ be a smooth real $(1,1)$-form near $\overline D$. Suppose $u,v\in C(\overline D)$ are $\vartheta$-plurisubharmonic,
\[
  (\vartheta+\ddc u)^n\geq F\beta^n,
  \qquad
  (\vartheta+\ddc v)^n\leq F\beta^n,
\]
where $F>0$ is continuous on $D$. If $u\leq v$ on $\partial D$, then $u\leq v$ in $D$.
\end{proposition}

\begin{proof}
\emph{Step 1: $v$ admits no strict smooth lower test.} Suppose, to the contrary, that a $C^2$ function $q$ touches $v$ from below at a point $p\in D$ and that, at $p$,
\[
  \vartheta+\ddc q>0,
  \qquad
  \det_\beta(\vartheta+\ddc q)>F.
\]
\Needspace{6\baselineskip}
Choose a sufficiently small $a>0$. After replacing $q$ by $q-a|z-p|^2$, which makes the contact strict, and restricting to a small ball $B_r(p)$, there are fixed $\lambda,\eta>0$ such that
\[
  A:=\vartheta+\ddc q\geq\lambda\beta,
  \qquad A^n\geq(F+\eta)\beta^n.
\]
To compare with closed currents, freeze the background at $p$: put $\kappa=\sup_{B_r}\lVert\vartheta-\vartheta(p)\rVert_\beta$ and $B_0=\vartheta(p)+\kappa\beta$. Then $S=B_0-\vartheta$ satisfies $0\leq S\leq2\kappa\beta$. For small $s>0$ the set $E=\{v<q+s\}$ is nonempty and relatively compact in $B_r$. If $Q_0$ is a quadratic potential for the constant form $B_0$, then
\[
  P=\ddc(v+Q_0)=\vartheta+\ddc v+S,
  \qquad
  Q=\ddc(q+s+Q_0)=A+S\geq\lambda\beta
\]
are closed positive Hessians. Applying \eqref{eq:interior-power-comparison}, expanding $P$, and using $S\leq2\kappa\beta\leq2\kappa\lambda^{-1}Q$, we obtain
\[
  \int_E P^n\leq\int_E F\beta^n+e_\kappa\int_E P^n,
  \qquad
  e_\kappa=\left(1+\frac{2\kappa}{\lambda}\right)^n-1.
\]
On the other hand, \eqref{eq:interior-power-comparison} with $j=n$ gives
\[
  \int_E P^n\geq\int_E Q^n\geq\int_E(F+\eta)\beta^n.
\]
If $M$ bounds $F$ on the ball, the two estimates give $(1-e_\kappa)\eta\leq e_\kappa M$. Since $\kappa\to0$, and hence $e_\kappa\to0$, as $r\downarrow0$, this is impossible for small $r$. Thus no such lower test exists.

\emph{Step 2: comparison.} Suppose that $\max_D(u-v)>0$. We regularize $u$ to produce a forbidden lower test for $v$. Let $\chi_\delta$ be a nonnegative smooth mollifier of mass one and radius $\delta>0$. For small $t>0$ and $C\geq\sup_D|z|^2$, set
\[
  u_t=u+t(|z|^2-C).
\]
Then $u_t\leq v$ on $\partial D$, the positive maximum persists, and $T_t=\vartheta+\ddc u_t\geq t\beta$. \Cref{lem:interior-mixed-background} gives
\[
  T_t^n\geq(F^{1/n}+t)^n\beta^n.
\]
Choose a relatively compact neighborhood $W$ of the positive maximum set. By \cref{lem:interior-convolution},
\[
  (T_t*\chi_\delta)^n
  \geq\bigl((F^{1/n}+t)*\chi_\delta\bigr)^n\beta^n.
\]
Since $T_t\geq t\beta$, the convolution satisfies $T_t*\chi_\delta\geq t\beta$. Moreover,
\[
  \vartheta+\ddc(u_t*\chi_\delta)
  =T_t*\chi_\delta+(\vartheta-\vartheta*\chi_\delta),
\]
so if
$\epsilon_\delta=\|\vartheta-\vartheta*\chi_\delta\|_{C^0(W)}$, then
\[
 \vartheta+\ddc(u_t*\chi_\delta)
 \geq(1-\epsilon_\delta/t)(T_t*\chi_\delta).
\]
For fixed $t>0$, uniform convergence of the convolved density and
$\epsilon_\delta\to0$ show that the left-hand side is positive and has
determinant strictly larger than $F$ on $W$ for small $\delta$. Since $u_t*\chi_\delta\to u_t$ uniformly on $W$, the difference $u_t*\chi_\delta-v$ still attains its maximum over $\overline W$ at an interior point of $W$ for small $\delta$. Subtracting this maximum from $u_t*\chi_\delta$ produces a smooth function touching $v$ from below with the properties excluded in Step~1, a contradiction.
\end{proof}

\subsection{Bounded traces and smoothness}\label{sec:interior-criterion}

We begin with the criterion used in the last step of the proof. A bounded trace of a positive current bounds all its coefficient measures, so a continuous solution with bounded trace is a strong solution with $W^{2,q}$ regularity for every finite $q$. The equation then bounds the eigenvalues from below, so it becomes uniformly elliptic. Since the solution is not yet known to be $C^2$, the Evans--Krylov theory is applied to it as a viscosity solution of a concave extension of the operator.

\begin{proposition}\label{prop:interior-criterion}
Let $B\subset\mathbb C^n$ be a coordinate ball, let $\vartheta$ be a smooth real $(1,1)$-form there and let $F\in C^\infty(B)$ be strictly positive. Suppose $u$ is continuous, $T=\vartheta+\ddc u\geq0$, and
\begin{equation}\label{eq:interior-local-ma}
  T^n=F\beta^n.
\end{equation}
For a point $p\in B$, the following conditions, each required on some neighborhood of $p$, are equivalent:
\begin{enumerate}
\item $u$ is smooth;
\item $T\geq\delta\beta$ for some $\delta>0$;
\item $\operatorname{tr}_\beta T$ has a locally bounded measure density.
\end{enumerate}
\end{proposition}

\begin{proof}
Choose $C$ so that $v=u+C|z|^2$ is plurisubharmonic, and use nonnegative radial mollifiers $\chi_\epsilon$ for which $v_\epsilon=v*\chi_\epsilon\downarrow v$. The smooth matrix convolutions $T_\epsilon=T*\chi_\epsilon$ can be written as
\[
  T_\epsilon=(\vartheta*\chi_\epsilon-C\beta)+\ddc v_\epsilon.
\]
Bedford--Taylor monotone continuity and smooth convergence of the background imply $T_\epsilon^k\to T^k$ for every $k$.

\emph{(2)$\Rightarrow$(3).} Assume (2). For every Hermitian matrix $M\geq\delta I$,
\[
  \det M\geq\delta^{n-1}\bigl(\operatorname{tr} M-(n-1)\delta\bigr).
\]
Apply this to the regularizations and pass to currents. Equation \eqref{eq:interior-local-ma} gives
\[
  \operatorname{tr}_\beta T\leq\delta^{1-n}F+(n-1)\delta,
\]
which is (3).

\emph{(3)$\Rightarrow$(1).} Assume (3), say $\operatorname{tr}_\beta T\leq L$. Positivity makes every coefficient of the matrix-valued current $T$ a bounded density. Hence the Euclidean Laplacian of $u$ is locally bounded, and interior Poisson estimates give $u\in W^{2,q}_{\mathrm{loc}}$ for every finite $q$. Mollification converges in $W^{2,q}$; taking $q>2n$ identifies the bounded-potential determinant with the ordinary almost-everywhere determinant. Thus $\det_\beta T=F$ almost everywhere. On a smaller ball with $F\geq m>0$, all eigenvalues of $T$ are at most $L$ and their product is at least $m$, so
\[
  \frac{m}{L^{n-1}}I\leq T\leq LI
\]
almost everywhere.

\Needspace{9\baselineskip}
It remains to show that this strong solution is smooth. Let $P$ be the real linear map from real symmetric $2n\times2n$ Hessian matrices to Hermitian $n\times n$ matrices, sending $D^2u$ to the coefficient matrix of $\ddc u$ relative to $\beta$, and let $\Theta(x)$ be the matrix of $\vartheta$. With $\delta=mL^{1-n}$, define the concave extension
\[
  G(x,M)=\inf_{\delta I\leq A\leq LI}
  \left\{\log\det A+\operatorname{tr}\!\left[A^{-1}\bigl(\Theta(x)+P(M)-A\bigr)\right]\right\}-\log F(x).
\]
For fixed $A$, write $\mathcal E_A(x,M)$ for the expression in braces
minus $\log F(x)$. The tangent-plane formula for $\log\det$ gives
\[
  G(x,D^2u)=0\quad\text{almost everywhere}.
\]
The operator $G$ is concave and uniformly elliptic on real symmetric matrices. The affine operators $\mathcal E_A$ have constant principal matrices with eigenvalues bounded above and below in terms of $\delta$ and $L$. Their $x$-dependence is smooth uniformly in $A$, so $G$ has a uniform local Lipschitz modulus in $x$.

We claim that $u$ is a viscosity solution of $G(x,D^2u)=0$. Suppose that a smooth function $\psi$ touches $u$ from above at $x_0$ and that $G(x_0,D^2\psi)<0$. Some fixed matrix $A$ then satisfies $\mathcal E_A(x,D^2\psi)<-\eta$ on a small ball. Since $G(x,D^2u)=0$ almost everywhere, $\mathcal E_A(x,D^2u)\geq0$ there. Thus $u-\psi$ has an interior maximum and satisfies
\[
  \operatorname{tr}\!\left(A^{-1}P(D^2(u-\psi))\right)\geq\eta
  \quad\text{almost everywhere}.
\]
After a small quadratic perturbation makes the contact strict, this contradicts the Aleksandrov maximum principle \cite[Chapter~9, Theorem~9.1]{GT01}. If instead $\psi$ touches $u$ from below and $G(x_0,D^2\psi)>0$, then $\mathcal E_A(x,D^2\psi)>\eta$ for every admissible $A$ on a small ball. At almost every point choose an admissible $A(x)$ with $\mathcal E_{A(x)}(x,D^2u)<\eta/2$. A measurable choice is obtained by taking the first qualifying element of a countable dense family. The function $\psi-u$ then violates the same maximum principle for the measurable uniformly elliptic coefficients $A(x)^{-1}$. This proves the claim.

The classical variable-coefficient Evans--Krylov theorem gives $u\in C^{2,\alpha}_{\mathrm{loc}}$ for some $0<\alpha<1$. The original equation is now uniformly elliptic with smooth data. Difference quotients satisfy linear equations whose principal coefficients are averages of inverse matrices along positive definite segments. Interior Schauder estimates first give $u\in C^{3,\alpha}_{\mathrm{loc}}$ and then, by iteration, $u\in C^\infty$. This proves (3)$\Rightarrow$(1).

\emph{(1)$\Rightarrow$(2).} If $u$ is smooth, then $T$ is a smooth semipositive form with determinant $F>0$, so it is positive definite; on a smaller neighborhood this gives (2).
\end{proof}

\Needspace{12\baselineskip}
\subsection{Smooth exact approximants}\label{sec:interior-approximants}

We now construct the approximants of the first step. The smooth Dirichlet theorem, \cref{lem:ball-smooth-dirichlet}, applies to every smooth positive density $F$. Approximating the H\"older boundary values of $\varphi$ by smooth functions therefore gives smooth solutions, and the point is to control them quantitatively. The estimates below bound the trace of a smooth solution by a power of the $C^4$ norm of its boundary value. The $C^4$ norms of the boundary approximations grow polynomially while their uniform errors decay polynomially, so both the convergence and the trace growth have polynomial rates.

In this subsection balls are centered at the origin. On the boundary sphere, $\nu$ is the outward Euclidean unit normal and $\nabla_{\tan}$ denotes tangential differentiation. For a positive form $T$, inner products $\langle\cdot,\cdot\rangle_T$ and norms $|\cdot|_T$ of $(1,0)$-covectors are computed with the inverse of the coefficient matrix of $T$.

\Needspace{20\baselineskip}
\begin{lemma}\label{lem:interior-trace-peak}
Let $B$ be a coordinate ball, let $\vartheta$ and $F>0$ be smooth near $\overline B$, and let $q$ be smooth with $R=\vartheta+\ddc q\geq\varepsilon\beta$ for a fixed $\varepsilon>0$. Fix a finite $L\geq0$. If $u\in C^\infty(B)$ satisfies
\[
  T=\vartheta+\ddc u>0,
  \qquad T^n=F\beta^n,
  \qquad |u-q|\leq L,
\]
put $S=\operatorname{tr}_\beta T$ and $v=u-q$. For $A>0$, define
\[
  Q_A=\log S-Av+\mathrm{e}^{-v}.
\]
There are $A_0,C$, depending only on the fixed data and $L$, such that at every interior local maximum of $Q_A$, with $A\geq A_0$,
\begin{equation}\label{eq:interior-trace-peak}
  S\leq C\sqrt A.
\end{equation}
The constants do not depend on any bound for the gradient of $u$.
\end{lemma}

\begin{proof}
This is the maximum-principle computation behind the classical second-order estimate. The non-closedness of $\vartheta$ produces a first-order error term, which is absorbed by the convex term $\mathrm{e}^{-v}$.

At the point in question make a constant unitary change of coordinates so that the matrix of $T$ is diagonal, with eigenvalues $\lambda_i$. Write
\[
  \Delta_T=\sum_i\lambda_i^{-1}\partial_i\partial_{\bar i},
  \qquad H=\operatorname{tr}_T\beta.
\]
Differentiating $\log\det T=\log F$ and commuting derivatives gives
\begin{equation}\label{eq:interior-trace-differentiation}
  \Delta_TS=\Delta_\beta\log F+P+E_0,
  \quad
  P=\sum_{k,i,j}\frac{|\partial_kT_{i\bar j}|^2}{\lambda_i\lambda_j},
  \quad |E_0|\leq C_0H.
\end{equation}
The non-closedness of $\vartheta$ enters through the bounded vector
\[
  E_p=\sum_k(\partial_p\vartheta_{k\bar k}-\partial_k\vartheta_{p\bar k}).
\]
\Needspace{6\baselineskip}
Since $\partial_pS-E_p=\sum_k\partial_kT_{p\bar k}$, weighted Cauchy--Schwarz gives $|\partial S-E|_T^2\leq SP$. From \eqref{eq:interior-trace-differentiation},
\begin{equation}\label{eq:interior-log-trace}
  \Delta_T\log S
  \geq-\frac{C_0(1+H)}S-
  \frac2S\operatorname{Re}\langle\partial\log S,E\rangle_T.
\end{equation}

Set $g_A(s)=-As+\mathrm{e}^{-s}$. At a local maximum of $Q_A$, $\partial\log S=-g_A'(v)\partial v$ and $\Delta_Tv=n-\operatorname{tr}_TR$. Combining these identities with \eqref{eq:interior-log-trace}, then completing the square in $\partial v$, yields
\begin{equation}\label{eq:interior-peak-inequality}
  0\geq
  \left[\varepsilon A-\frac{C_2}{S}-\frac{C_2(A+1)^2}{S^2}\right]H
  -C_3(A+1)-\frac{C_2}{S}.
\end{equation}
Here $|E|_T^2\leq CH$ and the bounds for $v$ have been used. If $S\geq K\sqrt A$ and $K$ is large, the bracket is at least $\varepsilon A/2$. Thus $H\leq C$ independently of $A$ and $u$. The determinant upper bound then gives $S\leq C'$. Enlarging $K$ makes this incompatible with $S\geq K\sqrt A$, proving \eqref{eq:interior-trace-peak}.
\end{proof}

We shall also use the following variant, in which a gradient bound replaces the convex term. If $|\nabla(u-q)|_\beta$ is bounded, then at an interior maximum of $\log S-A(u-q)$ the trace is bounded independently of $A$ for all large $A$. Indeed, at the maximum $\partial\log S=A\partial(u-q)$, and \eqref{eq:interior-log-trace} gives
\[
  0\geq\left[A\varepsilon-\frac CS-\frac{CA}{S}\right]H-nA-\frac CS.
\]
For $S$ larger than a fixed constant the bracket is at least $A\varepsilon/2$; the determinant again controls $S$.

\Needspace{19\baselineskip}
\paragraph{A polynomial boundary trace estimate}

\begin{proposition}\label{prop:interior-global-trace}
Let $B$ be the unit ball. Fix smooth $\vartheta$, smooth $F>0$ near $\overline B$, and $M<\infty$. There are $C,N<\infty$, depending only on these fixed data, such that every smooth solution $u$ with boundary value $g\in C^\infty(\partial B)$ of
\[
  T=\vartheta+\ddc u>0,
  \qquad T^n=F\beta^n,
  \qquad \lVert u\rVert_{L^\infty(B)}\leq M,
  \qquad u|_{\partial B}=g,
\]
satisfies
\begin{equation}\label{eq:interior-polynomial-global-trace}
  \sup_{\overline B}\operatorname{tr}_\beta T
  \leq C\bigl(1+\lVert g\rVert_{C^4(\partial B)}\bigr)^N.
\end{equation}
\end{proposition}

\begin{proof}
We follow the classical boundary estimates for the Dirichlet problem on a ball and keep track of the dependence on $L=1+\lVert g\rVert_{C^4}$. Barriers bound the first derivatives, and differentiation along rotations bounds the mixed normal--tangential second derivatives on the boundary. A lower bound for the tangential block of $T$ then bounds the double normal derivative through the equation, and \cref{lem:interior-trace-peak} reduces the interior bound to the boundary bound.

\emph{Step 1: first derivatives.} Fix a sufficiently large constant $K>0$, depending only on the fixed data. Extend $g$ to $G$ on $\overline B$ with $\lVert G\rVert_{C^4}\leq CL$. For $A=CL$ large enough,
\[
  \underline u=G+A(|z|^2-1),
  \qquad
  \underline T=\vartheta+\ddc\underline u\geq K\beta,
  \qquad
  \underline T^n\geq F\beta^n.
\]
\Cref{prop:interior-comparison} gives $\underline u\leq u$. If $h$ solves $\Delta_\beta h=-C_\vartheta$ with $h=g$ on $\partial B$ and $C_\vartheta\geq\sup\operatorname{tr}_\beta\vartheta$, then $u\leq h$. Boundary differentiation of these barriers gives
\begin{equation}\label{eq:interior-first-boundary}
  |\partial_\nu u|+|\nabla_{\tan}u|\leq CL
  \quad\text{on }\partial B.
\end{equation}

Let $H=\operatorname{tr}_T\beta$, let $\mathcal L=\operatorname{tr}_T\ddc$ be the linearized operator, and let $w=u-\underline u\geq0$ be the barrier function used below. Since $H\geq nF^{-1/n}$ is bounded below, taking $K$ fixed and sufficiently large gives
\begin{equation}\label{eq:interior-w-barrier}
  \mathcal Lw=n-\operatorname{tr}_T\underline T\leq-(1+H).
\end{equation}
Let $E=\sum_k(z_k\partial_{z_k}+\bar z_k\partial_{\bar z_k})$ be the real radial field. Write $\Theta$ for the coefficient matrix of $\vartheta$. Differentiating $\log\det T=\log F$ and commuting $E$ with complex second derivatives yields
\[
  \mathcal L(Eu)=E(\log F)+2n-\operatorname{tr}_T(E\Theta+2\Theta).
\]
Hence $|\mathcal L(Eu)|\leq C(1+H)$. The same estimate holds for $Zu$ when $Z$ is an infinitesimal unitary rotation. Applying the minimum principle to $CL+C'w\pm Yu$, with $Y=E$ or $Z$, gives
\begin{equation}\label{eq:interior-first-fields}
  |Eu|+|Zu|\leq CL\quad\text{on }B.
\end{equation}

\emph{Step 2: mixed second derivatives on the boundary.} For a rotation $Z$, the function $R_Z=Z(u-G)$ vanishes on $\partial B$ and satisfies $|\mathcal LR_Z|\leq CL(1+H)$. The barriers $CLw\pm R_Z$ and \eqref{eq:interior-w-barrier} imply $|\partial_\nu R_Z|\leq CL^2$ on the boundary. Rotations span the real tangent space of $\partial B$, so all real mixed normal--tangential second derivatives are $O(L^2)$. Pure tangential second derivatives are $O(L)$ by differentiating $u=g$ twice and using \eqref{eq:interior-first-boundary}. Thus all complex Hessian entries except the complex normal--normal entry are $O(L^2)$.

\emph{Step 3: the double normal derivative.} It remains to bound the tangential block of $T$ from below. For $n\geq2$, let
\[
  \lambda_0=\min\{T_p(\xi,\bar\xi):p\in\partial B,
  \ \xi\in T_p^{1,0}(\partial B),\ |\xi|_\beta=1\}.
\]
Choose a minimizing pair $(p,\xi)$ and extend $\xi$ to a smooth unit tangential field near $p$, with uniformly bounded derivatives. On the sphere, any smooth function vanishing there satisfies
\[
  \ddc(u-G)(\xi,\bar\xi)=\frac12E(u-G).
\]
Define
\[
  a(z)=(\vartheta+\ddc G)_z(\xi(z),\overline{\xi(z)})-\frac12EG(z),
  \qquad
  P=a+\frac12Eu-\lambda_0.
\]
Then $P\geq0$ on the nearby spherical boundary and $P(p)=0$.
The boundary first-derivative bound also gives $\lambda_0\leq CL$ by
the tangential identity above. Consequently
\[
  |P|\leq CL,
  \qquad
  |\mathcal LP|\leq CL(1+H).
\]
On a fixed cap $D=B\cap B_r(p)$ consider
\[
  \Psi=P+C_1Lw+C_2L|z-p|^2.
\]
Choose $C_2$ so that $\Psi\geq0$ on the interior spherical part of $\partial D$, and then $C_1$ so that $\mathcal L\Psi\leq0$ by \eqref{eq:interior-w-barrier}. The minimum principle yields $\Psi\geq0$. At $p$ the normal derivative of $|z-p|^2$ vanishes, and therefore
\[
  \partial_\nu P+C_1L\partial_\nu w\leq0.
\]
Equations \eqref{eq:interior-first-boundary}--\eqref{eq:interior-first-fields} now give $D^2u(\nu,\nu)(p)\leq CL^2$, hence $\operatorname{tr}_\beta T(p)\leq CL^2$. Positivity and $\det_\beta T=F$ imply
\begin{equation}\label{eq:interior-tangential-lower}
  \lambda_0\geq cL^{-2(n-1)}.
\end{equation}

\Needspace{5\baselineskip}
At any boundary point write
\[
  T=\begin{pmatrix}B_0&b\\ b^*&d\end{pmatrix}
\]
in tangential--normal blocks, with $B_0$ an $(n-1)\times(n-1)$ Hermitian matrix, $b\in\mathbb C^{n-1}$ and $d\in\mathbb R$. We have $|b|\leq CL^2$ and $B_0\geq\lambda_0I$. The determinant equation and the Schur complement give
\[
  d=\frac{F}{\det B_0}+b^*B_0^{-1}b
  \leq(\sup F)\lambda_0^{-(n-1)}+CL^4\lambda_0^{-1},
\]
which is polynomial in $L$. For $n=1$, the trace is $F$ and this argument is unnecessary.

\emph{Step 4: the interior.} Choose a fixed quadratic $q$ with $\vartheta+\ddc q\geq\varepsilon\beta$. Since $|u-q|\leq M+\sup_B|q|$, \cref{lem:interior-trace-peak} applies with this bound in place of its constant $L$; we use it with one fixed $A$. At an interior maximum of its auxiliary function the trace is uniformly bounded, while at a boundary maximum the polynomial boundary bound of Steps~2 and~3 applies. Since the remaining terms are uniformly bounded, \eqref{eq:interior-polynomial-global-trace} follows.
\end{proof}

\Needspace{23\baselineskip}
\paragraph{Exact approximation}

We apply these estimates to smooth approximations of the boundary values of $\varphi$ on a ball in $\Omega$. Recall that $\theta$, $\omega$, $c$ and $f$ are the data of \cref{thm:nonclosed-interior-smoothness}.

\Needspace{18\baselineskip}
\begin{proposition}\label{prop:interior-tame-approximants}
Let $B\Subset\Omega$ be a coordinate ball, identified with $\{|z|<r_0\}$, with closure in a larger coordinate chart, and put
\[
  F=\frac{cf\,\omega^n}{\beta^n}.
\]
There are $a>0$, $b,C<\infty$, and $u_j\in C^\infty(\overline B)$ such that
\begin{equation}\label{eq:interior-exact-approximants}
  \theta+\ddc u_j>0,
  \qquad
  (\theta+\ddc u_j)^n=F\beta^n,
\end{equation}
and
\begin{equation}\label{eq:interior-tame-rates}
  \lVert u_j-\varphi\rVert_{L^\infty(B)}\leq Cj^{-a},
  \qquad
  \sup_{\overline B}\operatorname{tr}_\beta(\theta+\ddc u_j)\leq Cj^b.
\end{equation}
\end{proposition}

\begin{proof}
Since $\varphi$ is H\"older continuous near $\overline B$, mollification at scale $1/j$ gives smooth boundary approximations $g_j$ with
\[
  \lVert g_j-\varphi\rVert_{L^\infty(\partial B)}\leq Cj^{-a},
  \qquad
  \lVert g_j\rVert_{C^4(\partial B)}\leq Cj^4.
\]
\Needspace{7\baselineskip}
Choose a fixed quadratic $q$ with $\eta=\theta+\ddc q>0$ near $\overline B$. Extend $g_j$ smoothly to $G_j$. For a sufficiently large, polynomially bounded $A_j$, the function
\[
  \underline w_j=G_j-q+A_j(|z|^2-r_0^2)
\]
has boundary value $g_j-q$ and satisfies
\[
  \eta+\ddc\underline w_j
  =\theta+\ddc G_j+A_j\beta>0,
  \qquad
  (\eta+\ddc\underline w_j)^n\geq F\beta^n.
\]
Apply the smooth Dirichlet theorem recalled at the beginning of this subsection with the background $\eta$, the density $F$ and the boundary value $g_j-q$. It gives a smooth solution $w_j$ of $(\eta+\ddc w_j)^n=F\beta^n$ with $\eta+\ddc w_j>0$; put $u_j=w_j+q$. Comparing $u_j$ with $\varphi\pm\lVert g_j-\varphi\rVert_\infty$ by \cref{prop:interior-comparison} proves the first bound in \eqref{eq:interior-tame-rates}. The functions $u_j$ are therefore uniformly bounded. \Cref{prop:interior-global-trace}, applied after rescaling $B$ to the unit ball, and the polynomial $C^4$ boundary norm give the second bound.
\end{proof}

\subsection{Localization near the nonregular set}\label{sec:interior-localization}

In the auxiliary function of \cref{lem:interior-trace-peak}, the term $-A(u_j-q)$ penalizes points where $u_j-q$ is large. A quadratic gap of $u-q$ along the nonregular set therefore suppresses the contribution of the boundary of a neighborhood of that set. Taking $A$ proportional to $\log j$ makes this suppression beat the polynomial growth of the boundary trace, and near the point of minimal $u-q$ the trace then grows more slowly than every positive power of $j$. Interpolation with the polynomial approximation rate gives uniform $C^{1,\alpha}$ bounds.

Distances and norms below are Euclidean. By a smooth exact solution for $(\vartheta,F)$ we mean a smooth function $u_j$ with $\vartheta+\ddc u_j>0$ and $(\vartheta+\ddc u_j)^n=F\beta^n$.

\Needspace{28\baselineskip}
\begin{proposition}\label{prop:interior-subpower}\label{prop:interior-uniform-c1a}
For $R>0$, let $B=B_R(0)$, let $\vartheta$ and $F>0$ be smooth near $\overline B$, and let $u\in C(\overline B)$. Suppose $q$ is smooth and $\vartheta+\ddc q\geq\varepsilon\beta$ for some $\varepsilon>0$. Put $v=u-q$ and $m=v(0)$, and fix $\kappa>0$. Let $\mathcal S\subset\overline B$ be closed, contain $0$, and satisfy
\begin{equation}\label{eq:interior-quadratic-separation}
  v(0)=m,
  \qquad
  v(z)\geq m+\kappa|z|^2\quad(z\in\mathcal S).
\end{equation}
Assume that smooth exact solutions $u_j$, $j\geq2$, satisfy, for fixed constants $a>0$, $b\geq0$ and $C>0$,
\begin{align}
  \lVert u_j-u\rVert_\infty&\leq Cj^{-a}, \label{eq:interior-localization-a}\\
  \sup_{\overline B}\operatorname{tr}_\beta(\vartheta+\ddc u_j)&\leq Cj^b, \label{eq:interior-localization-b}\\
  \sup_j\sup_K\operatorname{tr}_\beta(\vartheta+\ddc u_j)&<\infty
  \quad\text{for every }K\Subset\overline B\setminus\mathcal S. \label{eq:interior-localization-c}
\end{align}
Then for every $\gamma>0$ there are a neighborhood $U_\gamma$ of $0$ and $C_\gamma$ such that
\begin{equation}\label{eq:interior-subpower-trace}
  \operatorname{tr}_\beta(\vartheta+\ddc u_j)\leq C_\gamma j^\gamma
  \quad\text{on }U_\gamma.
\end{equation}
For every $0<\alpha<1$, there is also a neighborhood $U$ of $0$ on which $u_j\to u$ in $C^{1,\alpha}$. In particular the approximants are uniformly bounded in that norm and $u\in C^{1,\alpha}(U)$.
\end{proposition}

\begin{proof}
\emph{Step 1: the trace bound \eqref{eq:interior-subpower-trace}.} Write $e_j=Cj^{-a}$, fix $d=\kappa R^2/2$, and choose $D>(b+1)/d$. Choose $\eta>0$ so that $2D\eta<\gamma/2$. For small $\delta>0$, the neighborhood
\[
  W_\delta=\{z\in B:\operatorname{dist}(z,\mathcal S)<\delta\}
\]
contains $0$, has $v\geq m-\eta$ on its closure, and satisfies $v\geq m+d$ on $\partial B\cap\overline W_\delta$ with the fixed $d$ above. These properties follow from \eqref{eq:interior-quadratic-separation}, continuity, and compactness. The remaining part of the boundary of $W_\delta$ is a compact subset of $\overline B\setminus\mathcal S$, where \eqref{eq:interior-localization-c} applies.

Set $A_j=D\log j$. Apply the maximum principle on $\overline W_\delta$ to
\[
  Q_j=\log\operatorname{tr}_\beta(\vartheta+\ddc u_j)-A_j(u_j-q)+\mathrm{e}^{-(u_j-q)}.
\]
At an interior maximum, \cref{lem:interior-trace-peak} gives $\log\operatorname{tr}_\beta(\vartheta+\ddc u_j)\leq C+\frac12\log A_j$. On the part of $\partial W_\delta$ inside $B$, which stays away from $\mathcal S$, the trace has a bound independent of $j$. On the outer part $\partial B\cap\overline W_\delta$, \eqref{eq:interior-localization-b} and the gap $d$ give
\[
  Q_j\leq C+b\log j-A_j(m+d-e_j),
\]
and the choice $Dd>b+1$ makes this smaller than the bound below for large $j$. Hence throughout $W_\delta$,
\[
  Q_j\leq C_\delta+\tfrac12\log A_j-A_j(m-\eta-e_j).
\]
On a small ball about $0$ where $v\leq m+\eta$, this implies
\[
  \log\operatorname{tr}_\beta(\vartheta+\ddc u_j)
  \leq C_\delta+\tfrac12\log A_j+A_j(2\eta+2e_j).
\]
By the choice of $\eta$, the logarithmic factor and $A_je_j$ are absorbed into the remaining power. This proves \eqref{eq:interior-subpower-trace}.

\emph{Step 2: $C^{1,\alpha}$ convergence.} On a smaller ball, the real Laplacian of $u_j$ is bounded by $Cj^\gamma$. Interior Poisson estimates give a $W^{2,q_0}$ bound for every finite $q_0$, and Sobolev embedding with $q_0$ large then gives, for every $\sigma<1$,
\[
  \lVert u_j\rVert_{W^{2,q_0}}+\lVert u_j\rVert_{C^{1,\sigma}}
  \leq Cj^\gamma.
\]
For $d_k=u_{2^{k+1}}-u_{2^k}$,
\[
  \lVert d_k\rVert_\infty\leq C2^{-ka},
  \qquad
  \lVert d_k\rVert_{C^{1,\sigma}}\leq C2^{k\gamma}.
\]
If $\alpha<\sigma$ and $t=(1+\alpha)/(1+\sigma)$, H\"older interpolation gives
\[
  \lVert d_k\rVert_{C^{1,\alpha}}
  \leq C2^{k[-a(1-t)+\gamma t]}.
\]
Since $\gamma$ may be arbitrarily small, the exponent can be made negative. The dyadic series converges in $C^{1,\alpha}$ to its uniform limit $u$. If $2^k\leq j<2^{k+1}$, interpolation applied to $u_j-u_{2^k}$ gives the same decaying bound. Hence the full sequence converges to $u$ in $C^{1,\alpha}$ on a smaller ball.
\end{proof}

\begin{proposition}\label{prop:interior-regular-stability}
Suppose $u_j$ solve the same smooth equation as in \cref{prop:interior-subpower}, satisfy \eqref{eq:interior-localization-a}--\eqref{eq:interior-localization-b}, and $u$ is smooth with $\vartheta+\ddc u>0$ on an open set $O\subset B$. Then, for every $K\Subset O$,
\begin{equation}\label{eq:interior-regular-stability}
  \sup_j\sup_K\operatorname{tr}_\beta(\vartheta+\ddc u_j)<\infty.
\end{equation}
\end{proposition}

\begin{proof}
Where $u$ is smooth and strictly $\vartheta$-plurisubharmonic, it provides its own quadratic separation, so the argument of \cref{prop:interior-subpower} applies around every point of $O$. A first application gives trace growth of order $\sqrt{\log j}$; the resulting gradient convergence then allows the variant of \cref{lem:interior-trace-peak} with a gradient bound, which removes this growth.

For $p$ in a fixed compact subset of $O$, take a uniformly sized ball $B_r(p)\Subset O$ and set
\[
  q_p(z)=u(z)-\kappa|z-p|^2
\]
with $\kappa>0$ small, so $\vartheta+\ddc q_p\geq\varepsilon\beta$. If $e_j=Cj^{-a}$ and $v_{j,p}=u_j-q_p$, then
\[
  v_{j,p}\geq-e_j,
  \quad v_{j,p}(p)\leq e_j,
  \quad v_{j,p}\geq\kappa r^2-e_j\text{ on }\partial B_r(p).
\]
With $A_j=D\log j$ and $D\kappa r^2>b+1$, the maximum argument of Step~1 in the proof of \cref{prop:interior-subpower} gives
\begin{equation}\label{eq:interior-sqrtlog-trace}
  \operatorname{tr}_\beta(\vartheta+\ddc u_j)(p)\leq C\sqrt{\log j}.
\end{equation}
The constants are uniform on compact subsets of $O$.

Put $T_j=\vartheta+\ddc u_j$ and $w_j=u_j-u$. On slightly larger compact sets, \eqref{eq:interior-sqrtlog-trace} bounds $|\Delta w_j|$ by $C\sqrt{\log j}$. Poisson estimates and Sobolev embedding give
\[
  \lVert w_j\rVert_{C^{1,\sigma}}\leq C_\sigma\sqrt{\log j},
  \qquad \sigma<1.
\]
Interpolating with $\lVert w_j\rVert_\infty\leq Cj^{-a}$ shows $\nabla w_j\to0$ uniformly on smaller compact sets. The gradients of $v_{j,p}$ are therefore uniformly bounded. Repeat the maximum argument with $\log\operatorname{tr}_\beta T_j-A_jv_{j,p}$, now using the variant of \cref{lem:interior-trace-peak} with a gradient bound. Its interior trace bound is now independent of $A_j$, while the boundary term is suppressed as before. Evaluation at $p$ gives
\[
  \operatorname{tr}_\beta T_j(p)\leq C\exp(2A_je_j)\leq C,
\]
proving \eqref{eq:interior-regular-stability}.
\end{proof}



\Needspace{24\baselineskip}
\subsection{Derivative estimates and bounded traces}\label{sec:interior-calabi}

Uniform $C^{1,\alpha}$ bounds control averages of the trace $\operatorname{tr}_\beta(\vartheta+\ddc u_j)=\operatorname{tr}_\beta\vartheta+\Delta_\beta u_j$ over small balls. To pass from averages to pointwise values we need a bound for the gradient of the trace, that is, a third-order estimate, whose dependence on the ellipticity constants is polynomial. We prove it by Calabi's maximum-principle argument, computed with the Chern connection of the solution so that the torsion of the nonclosed background appears explicitly. The resulting averaging inequality is nonlinear in the trace, and it closes under iteration when $\alpha$ is sufficiently close to one.

In this subsection $B$ is the unit ball and $B_r(p)$ is the Euclidean ball of radius $r$ centered at $p$. For a Hermitian metric $g$, $\partial g$ denotes the ordinary coordinate derivatives of its coefficient matrix, while tensor derivatives in the proofs are taken with respect to the Chern connection of $g$.

\begin{proposition}\label{prop:interior-calabi}
Fix smooth $\vartheta$ and smooth $F>0$ near the closed unit ball. There are $C,N<\infty$, depending only on finitely many derivatives of the fixed data, such that any smooth solution
\[
  g=\vartheta+\ddc u,
  \qquad g^n=F\beta^n,
  \qquad M^{-1}\beta\leq g\leq M\beta,
  \qquad M\geq1,
\]
satisfies
\begin{equation}\label{eq:interior-calabi-polynomial}
  \sup_{B_{1/2}}|\partial g|_\beta\leq CM^N.
\end{equation}
If only $\sup_B\operatorname{tr}_\beta g\leq L$ is given, the same bound holds with a fixed power of $1+L$. On a ball of radius $r\leq1$,
\begin{equation}\label{eq:interior-gradient-trace-scaled}
  \sup_{B_{r/2}}|\nabla\operatorname{tr}_\beta g|_\beta
  \leq Cr^{-1}(1+L)^N.
\end{equation}
\end{proposition}

\begin{proof}
The quantity $\mathcal R$ below is the squared norm of the difference between the Chern connection of $g$ and the flat connection. We derive a differential inequality for $\mathcal R$ in which the torsion enters through fixed derivatives of $\vartheta$, and combine it with the trace in a maximum principle; every constant is kept a power of $M$.

We compare the flat connection of $\beta$ with the Chern connection $\nabla$ of $g$. Write $g^{k\bar\ell}$ for the coefficients of the inverse matrix and $\bar\nabla$ for covariant differentiation in $(0,1)$ directions; tensor norms and inner products without subscript are computed with $g$. The difference of the two connections and its squared norm are
\[
  \Psi^k_{ij}=g^{k\bar\ell}\partial_i g_{j\bar\ell},
  \qquad \mathcal R=|\Psi|_g^2.
\]
The torsion of $\nabla$ is determined by the fixed background, because the third derivatives of $u$ cancel in it. Its lowered and raised components are
\begin{equation}\label{eq:interior-torsion-background}
  C_{ij\bar\ell}
  =\partial_i\vartheta_{j\bar\ell}-\partial_j\vartheta_{i\bar\ell},
  \qquad
  \mathcal T^k_{ij}=g^{k\bar\ell}C_{ij\bar\ell}.
\end{equation}
\Needspace{8\baselineskip}
Consequently, for a fixed integer $N$,
\begin{align}
  |\mathcal T|&\leq CM^N,\nonumber\\
  |\nabla\mathcal T|+|\bar\nabla\mathcal T|
  &\leq CM^N(1+\sqrt{\mathcal R}),\label{eq:interior-torsion-derivatives}\\
  |\nabla^2\mathcal T|+|\nabla\bar\nabla\mathcal T|+|\bar\nabla^2\mathcal T|
  &\leq CM^N(1+\mathcal R+|\nabla\Psi|+|\bar\nabla\Psi|).\nonumber
\end{align}
Every factor here is either a fixed derivative of $\vartheta$, a connection factor $\Psi$, or a factor of $g^{\pm1}$; this explains the polynomial dependence.

For tensor fields, use $\Delta_g=g^{p\bar q}\nabla_p\nabla_{\bar q}$.
We need the following curvature identities. With
\[
  \operatorname{Rm}^k{}_{i\bar j l}=-\partial_{\bar j}\Psi^k_{il},
  \qquad
  \operatorname{Rm}_{i\bar j l\bar k}=g_{m\bar k}\operatorname{Rm}^m{}_{i\bar j l},
\]
set
\[
  \operatorname{Ric}^{(2)}_{l\bar k}=g^{i\bar j}\operatorname{Rm}_{i\bar j l\bar k},
  \qquad
  \operatorname{Ric}^{(1)}_{l\bar k}=g^{i\bar j}\operatorname{Rm}_{l\bar k i\bar j}
  =-\partial_l\partial_{\bar k}\log F.
\]
The two curvature interchanges are
\begin{align*}
  \operatorname{Rm}_{i\bar j l\bar k}-\operatorname{Rm}_{l\bar j i\bar k}
  &=-\nabla_{\bar j}C_{il\bar k},\\
  \operatorname{Rm}_{i\bar j l\bar k}-\operatorname{Rm}_{i\bar k l\bar j}
  &=-\nabla_i\overline{C}_{\bar j\bar k l}.
\end{align*}
They follow by differentiating the torsion identity and its conjugate. Successive interchange and contraction give
\begin{equation}\label{eq:interior-ricci-contraction}
  \operatorname{Ric}^{(2)}_{l\bar k}
  =\operatorname{Ric}^{(1)}_{l\bar k}
  -g^{i\bar j}\bigl(
    \nabla_l\overline C_{\bar j\bar k i}
    +\nabla_{\bar j}C_{il\bar k}
  \bigr).
\end{equation}
Metric compatibility, \eqref{eq:interior-torsion-derivatives}, and the fixed smooth formula for $\operatorname{Ric}^{(1)}$ imply
\begin{equation}\label{eq:interior-ricci2}
  |\operatorname{Ric}^{(2)}|\leq CM^N(1+\sqrt{\mathcal R}),
  \qquad
  |\nabla\operatorname{Ric}^{(2)}|
  \leq CM^N(1+\mathcal R+|\nabla\Psi|+|\bar\nabla\Psi|).
\end{equation}
Moreover,
\[
  \nabla_{\bar q}\Psi^k_{ij}=-\operatorname{Rm}^k{}_{i\bar q j},
\]
and the Chern--Bianchi identity, including its torsion term, is
\[
  \nabla_p\operatorname{Rm}^k{}_{i\bar q j}-\nabla_i\operatorname{Rm}^k{}_{p\bar q j}
  =-\mathcal T^r_{pi}\operatorname{Rm}^k{}_{r\bar q j}.
\]
Consequently,
\begin{equation}\label{eq:interior-laplacian-connection-difference}
  \Delta_g\Psi^k_{ij}
  =-\nabla_i(\operatorname{Ric}^{(2)})^k{}_j
  +g^{p\bar q}\mathcal T^r_{pi}\operatorname{Rm}^k{}_{r\bar q j}.
\end{equation}
If $A^k{}_l=(\operatorname{Ric}^{(2)})^k{}_l$, define its action on a tensor of type $(1,2)$ by
\[
  (A\mathbin{\cdot}\Psi)^k_{ij}
  =A^k{}_l\Psi^l_{ij}-A^l{}_i\Psi^k_{lj}-A^l{}_j\Psi^k_{il}.
\]
Differentiating the squared norm and using metric compatibility gives
\begin{equation}\label{eq:interior-laplacian-calabi-norm}
  \Delta_g\mathcal R
  =|\nabla\Psi|^2+|\bar\nabla\Psi|^2
   +2\operatorname{Re}\langle\Delta_g\Psi,\Psi\rangle
   -\operatorname{Re}\langle A\mathbin{\cdot}\Psi,\Psi\rangle.
\end{equation}
There is no mixed torsion term because the Chern torsion has type $(2,0)$. Using \eqref{eq:interior-torsion-derivatives}, \eqref{eq:interior-ricci2}, and \eqref{eq:interior-laplacian-connection-difference} in \eqref{eq:interior-laplacian-calabi-norm}, Young's inequality absorbs the products of $\sqrt{\mathcal R}$ with $|\nabla\Psi|+|\bar\nabla\Psi|$. We obtain
\begin{equation}\label{eq:interior-calabi-differential}
  \Delta_g\mathcal R\geq-CM^N(\mathcal R^{3/2}+1).
\end{equation}
Let $s=\operatorname{tr}_\beta g$. Direct differentiation as in \eqref{eq:interior-trace-differentiation} gives
\begin{equation}\label{eq:interior-trace-calabi-coupling}
  \Delta_gs\geq M^{-1}\mathcal R-CM,
  \qquad
  |\partial s|_g^2\leq nM^2\mathcal R.
\end{equation}
Take a cutoff $\zeta$ supported in $B$, equal to one on $B_{1/2}$, set $\eta=\zeta^2$ and $h=2nM-s$, so $nM\leq h\leq2nM$, and maximize
\[
  Q=\frac{\eta\mathcal R}{h}.
\]
At a positive interior maximum of $Q$, the identity
$\partial\mathcal R/\mathcal R=\partial h/h-\partial\eta/\eta$ gives
\[
 \begin{split}
 0\geq\frac{\Delta_gQ}{Q}
 &=\frac{\Delta_g\mathcal R}{\mathcal R}
   +\frac{\Delta_g\eta}{\eta}-\frac{\Delta_gh}{h}
   -\frac{2|\partial\eta|_g^2}{\eta^2}
   +\frac{2\operatorname{Re}\langle\partial\eta,\partial h\rangle_g}{\eta h}.
 \end{split}
\]
We have $|\partial\eta|_g^2\leq CM\eta$,
$|\Delta_g\eta|\leq CM$ and
$|\partial h|_g\leq\sqrt n M\sqrt{\mathcal R}$.
Insert \eqref{eq:interior-calabi-differential} and
\eqref{eq:interior-trace-calabi-coupling} and multiply by $\eta Q$.
The positive term is $M^{-1}Q^2$. The two terms involving
$\sqrt{\mathcal R}$ are bounded in absolute value by a fixed power of
$M$ times $Q^{3/2}$, since $\mathcal R=hQ/\eta$ and $h\leq2nM$.
The remaining cutoff terms are bounded by $CM^N(Q+1)$. Therefore
\[
 0\geq M^{-1}Q^2-CM^N(Q^{3/2}+Q+1).
\]
Thus $Q$, and then $\mathcal R$ on $B_{1/2}$, are bounded by a fixed polynomial in $M$. Since $|\partial g|_\beta^2\leq M^3\mathcal R$, this proves \eqref{eq:interior-calabi-polynomial}.

If $s\leq L$, the determinant lower bound gives $\lambda_{\min}(g)\geq(\inf F)L^{1-n}$, so $M$ is a fixed power of $1+L$. For a ball $B_r(p)$, set $u_r(z)=r^{-2}u(p+rz)$ and use the
coefficient matrices $\vartheta_r(z)=\vartheta(p+rz)$,
$g_r(z)=g(p+rz)$ and $F_r(z)=F(p+rz)$. Then
$g_r=\vartheta_r+\ddc u_r$ in the unit coordinates and
$\det g_r=F_r$. For $r\leq1$ the derivatives of the fixed data remain
uniformly bounded. The unit-ball estimate for the ordinary coefficient derivatives $\partial g_r$ and their conjugates gives the
factor $r^{-1}$ in \eqref{eq:interior-gradient-trace-scaled}.
\end{proof}

\Needspace{12\baselineskip}
\begin{proposition}\label{prop:interior-c1a-to-trace}
There is $\alpha_*<1$, depending only on $n$ and the exponent $N$ in \cref{prop:interior-calabi}, with the following property. Let $B_0$ be a Euclidean ball, fix a smooth real $(1,1)$-form $\vartheta$ and a smooth function $F>0$ near $\overline{B_0}$, and let $u_j$ be smooth exact solutions of $(\vartheta+\ddc u_j)^n=F\beta^n$ with $\vartheta+\ddc u_j>0$. If
\[
  \sup_j\lVert u_j\rVert_{C^{1,\alpha}(B_0)}<\infty
  \qquad\text{for some }\alpha_*<\alpha<1,
\]
then $\operatorname{tr}_\beta(\vartheta+\ddc u_j)$ is uniformly bounded on every smaller ball.
\end{proposition}

\begin{proof}
Fix a center for concentric balls $B_r\Subset B_0$, put $s_j=\operatorname{tr}_\beta(\vartheta+\ddc u_j)$, and set
\[
  M_j(r)=1+\sup_{B_r}s_j.
\]
For concentric balls $B_r\Subset B_R\Subset B_0$, $d=R-r\leq1$, \cref{prop:interior-calabi} gives
\begin{equation}\label{eq:interior-trace-gradient-M}
  |\nabla s_j|\leq Cd^{-1}M_j(R)^N
  \quad\text{on }B_{r+d/2}.
\end{equation}
Fix $p\in B_r$ and let $\chi_h$ be a nonnegative mollifier of radius $h\leq d/4$. Since
\[
  s_j=\operatorname{tr}_\beta\vartheta+\Delta_\beta u_j,
\]
integration by parts, after subtracting the constant vector $\nabla u_j(p)$, and the uniform $C^{1,\alpha}$ bound give
\[
  (s_j*\chi_h)(p)\leq C(1+h^{\alpha-1}).
\]
Equation \eqref{eq:interior-trace-gradient-M} gives
\[
  |s_j(p)-(s_j*\chi_h)(p)|
  \leq Cd^{-1}M_j(R)^Nh.
\]
Choose
\[
  h=\frac d4M_j(R)^{-N/(2-\alpha)}.
\]
Taking the supremum over $p\in B_r$ and absorbing fixed powers of $d^{-1}$, we obtain
\begin{equation}\label{eq:interior-nonlinear-iteration}
  M_j(r)\leq Cd^{-1}M_j(R)^q,
  \qquad
  q=\frac{N(1-\alpha)}{2-\alpha}.
\end{equation}
Choose $\alpha_*$ so close to one that $q<1$ for $\alpha>\alpha_*$. Choose $0<r_0<R_0$ with $B_{R_0}\Subset B_0$. Let $r_k\uparrow R_0$ with $r_{k+1}-r_k=(R_0-r_0)2^{-k-1}$. Iterating \eqref{eq:interior-nonlinear-iteration} gives
\[
  M_j(r_0)
  \leq
  \prod_{k=0}^{m-1}\left[C(r_{k+1}-r_k)^{-1}\right]^{q^k}
  M_j(r_m)^{q^m}.
\]
The infinite product converges because $\sum q^k(1+k)<\infty$. For each fixed $j$, the final factor tends to one as $m\to\infty$. The resulting bound is independent of $j$.
\end{proof}

\Needspace{8\baselineskip}
\subsection{Proof of the main theorem}\label{sec:interior-main-proof}

We now assemble the four steps outlined after \cref{thm:nonclosed-interior-smoothness}. The barrier enters only once, but globally: it locates a point of the nonsmooth set at which the quadratic separation required in \cref{prop:interior-subpower} holds.

\begin{proof}[of \cref{thm:nonclosed-interior-smoothness}]
If $n=1$, the equation is the linear identity
\[
  \ddc\varphi=cf\,\omega-\theta
  \quad\text{on }\Omega.
\]
Its right-hand side is smooth, so interior elliptic regularity gives $\varphi\in C^\infty(\Omega)$, and the original equation gives $\theta+\ddc\varphi=cf\,\omega>0$. We henceforth assume $n\geq2$.

Let
\[
  \mathcal S=\{p\in\Omega:\varphi\text{ is not smooth on any neighborhood of }p\}.
\]
This is closed in $\Omega$. Suppose $\mathcal S\neq\varnothing$. By \cref{lem:interior-proper-gap}, the restriction of $\varphi-\rho$ to $\mathcal S$ attains its minimum $m$ at some $p\in\mathcal S$.

Choose nested coordinate balls
\[
  B\Subset B_0\Subset\Omega
\]
centered at $p$. On $B_0$, \cref{prop:interior-tame-approximants} gives smooth exact Dirichlet solutions $u_j$ converging uniformly to $\varphi$ at a polynomial rate and having polynomial global trace growth.

On the open set $\Omega\setminus\mathcal S$ the function $\varphi$ is smooth, and $\theta+\ddc\varphi$ is semipositive with positive determinant, hence positive definite. \Cref{prop:interior-regular-stability} therefore gives a uniform trace bound for the $u_j$ on every compact subset of $B_0\setminus\mathcal S$.

In coordinates centered at $p$, take $\kappa>0$ small and set
\[
  q(z)=\rho(z)-\kappa|z|^2.
\]
After shrinking $B$, the lower bound $\theta+\ddc\rho\geq\omega$ gives $\theta+\ddc q>0$ on $\overline B$. If $v=\varphi-q$, then the minimizing property of $p$ gives
\[
  v(p)=m,
  \qquad
  v(z)=(\varphi-\rho)(z)+\kappa|z|^2
  \geq m+\kappa|z|^2
  \quad(z\in\mathcal S\cap\overline B).
\]
Thus \cref{prop:interior-subpower} applies on $B$ with the closed set $\mathcal S\cap\overline B$: the approximation rate and global polynomial trace bound come from \cref{prop:interior-tame-approximants}, and the bound away from $\mathcal S$ comes from \cref{prop:interior-regular-stability}. Its second assertion gives, for every $0<\alpha<1$, a neighborhood of $p$ on which the approximants have a uniform $C^{1,\alpha}$ bound.

Choose $\alpha>\alpha_*$ as in \cref{prop:interior-c1a-to-trace}, and apply that proposition on a ball centered at $p$ inside this neighborhood. It yields
\[
  \sup_j\operatorname{tr}_\beta(\theta+\ddc u_j)<\infty
\]
on a still smaller neighborhood of $p$. Since $u_j\to\varphi$ uniformly, $\ddc u_j\to\ddc\varphi$ as currents. The trace bound passes to the limit, so $\operatorname{tr}_\beta(\theta+\ddc\varphi)$ has a bounded density near $p$. \Cref{prop:interior-criterion} now implies that $\varphi$ is smooth near $p$, contradicting $p\in\mathcal S$.

Hence $\mathcal S=\varnothing$ and $\varphi\in C^\infty(\Omega)$. The form $\theta+\ddc\varphi$ is smooth and semipositive there, while
\[
  (\theta+\ddc\varphi)^n=cf\,\omega^n>0.
\]
Its nonnegative eigenvalues have positive product, so every eigenvalue is positive. Therefore $\theta+\ddc\varphi>0$ on $\Omega$.
\end{proof}



\begin{corollary}\label{cor:big-lp-interior-smoothness}
Let $0\leq f\in L^p(X,\omega^n)$, $p>1$, be smooth and strictly positive
on $\Omega$. Every minimal-singularity solution of
$(\theta+\ddc\varphi)^n=cf\,\omega^n$, $c>0$, is smooth and strictly
$\theta$-plurisubharmonic on $\Omega$.
\end{corollary}
\begin{proof}
The analytic complement of $\Omega$ has zero mass for both measures,
so it suffices to know the equation in the bounded-potential sense on
$\Omega$. Apply \cref{thm:big-minimal-holder} with density $cf$ to obtain
local H\"older continuity. All hypotheses of
\cref{thm:nonclosed-interior-smoothness} are now satisfied.
\end{proof}

By \cref{thm:big-minimal-normalized}, such a solution exists, with a uniquely determined $c>0$. For a smooth closed big representative on a compact K\"ahler manifold,
\cref{thm:big-class-ample-smoothness} applies this conclusion to the
classical full-mass solution and identifies the regular open set with
the ample locus.
